\documentclass[10pt,a4paper]{amsart}
\usepackage[margin=19mm]{geometry}
\usepackage{amsmath,amssymb,mathtools}
\usepackage[scr=rsfs,cal=euler]{mathalfa}
\usepackage{xfrac}
\usepackage[unicode,hidelinks,bookmarks=false]{hyperref}
\newcommand{\ie}{i.e.\ }
\newcommand{\cf}{cf.\ }
\newcommand{\resp}{respectively\ }
\newcommand{\Subsection}{\subsection}
\providecommand{\nobreakdash}{\nobreak\hbox{-}}
\setlength{\emergencystretch}{2em}
\multlinegap0pt
\usepackage{bm}
\usepackage{bbm}
\usepackage{dsfont}
\usepackage{xspace}
\usepackage{cancel}
\usepackage{mathtools}
\usepackage{amsthm}
\makeatletter
\@ifundefined{assumption}{\newtheorem{assumption}{Assumption}}{}
\@ifundefined{corollary}{\newtheorem{corollary}{Corollary}}{}
\@ifundefined{theorem}{\newtheorem{theorem}{Theorem}}{}
\makeatother
\DeclareMathOperator{\Tr}{Tr}
\title[On Regularization via Early Stopping for Least Squares Regre]{On Regularization via Early Stopping for Least Squares Regression}
\author{Rishi Sonthalia, Jackie Lok, Elizaveta Rebrova}
\date{Source: arXiv:2406.04425v2 (CC BY 4.0)}
\begin{document}
\maketitle

\noindent\emph{Source and licence.} On Regularization via Early Stopping for Least Squares Regression. Version arXiv:2406.04425v2 (CC BY 4.0), distributed under the Creative Commons Attribution 4.0 International licence, as recorded in the arXiv licence metadata for this exact version and shown on the abstract page https://arxiv.org/abs/2406.04425v2. Full rights notice: licenses/arxiv-2406.04425-CC-BY-4.0.txt.\par\smallskip

\noindent\textit{Editorial scope.} Counted unit: the Theorem whose source label is \texttt{thm:earlystopping} in the article (source0.tex lines 561--573, source label \texttt{thm:earlystopping}), delivered together with the article's own complete original proof of it (source0.tex lines 575--618, one \texttt{proof} environment of 44 lines). No mathematics is written, repaired or re-derived: statement and proof are the article's own text line for line. Required context retained uncounted: eq:ls-ridge (lines 152--154); assumption:noise (lines 500--502); cor:gen-risk-betak (lines 518--523); assumption:differentiable-extension (lines 542--544); assumption:spherical (lines 551--553). Every definition, display and label of those lines is retained unchanged. Omitted from this package and not redistributed: 1--151, 155--499, 503--517, 524--541, 545--550, 554--560, 574--574, 619--1276. Nothing inside a retained excerpt is omitted or altered: every retained body is delivered exactly as the article's lines are written, so no omission receipt of any kind accompanies this package. Citations: the retained excerpts carry no citation command at all, so no bibliography entry of the article is transcribed and no reference is redistributed. Reference adaptation: none was needed and none was made; every reference inside the delivered unit resolves inside it, so no unresolved reference or citation is produced. Printed slips kept literal: no printed word, symbol, hypothesis or number of the counted statement or of its proof was altered. Layout and class: the article's own class is \texttt{article} with packages including geometry, amsmath, amssymb, amsthm, tcolorbox, wrapfig, blkarray, booktabs, multirow, thm-restate, thmtools, subcaption, bbm, inputenc; the wrapper is the lane's pinned \texttt{amsart} editorial wrapper at 10pt on A4, which restates the theorem-like environments the retained text needs and adds the article's own macro definitions that the retained text needs (\texttt{\textbackslash Tr}), with extra wrapper packages (bm, bbm, dsfont, xspace, cancel, mathtools). The delivered numbering is the unit's own. Assets: no retained span contains a figure environment, a graphics-inclusion command, a PDF-page inclusion, a table or a TikZ picture; no image file is copied into this package, the package declares no asset dependency and no image file is redistributed with it. Licence: version arXiv:2406.04425v2 of this article is distributed under the Creative Commons Attribution 4.0 International licence, as recorded in the arXiv licence metadata for this exact version and shown on the abstract page https://arxiv.org/abs/2406.04425v2. A full rights notice is delivered at licenses/arxiv-2406.04425-CC-BY-4.0.txt. Characters: every retained excerpt is pure ASCII, so the delivered engine, XeLaTeX with its default Latin Modern text font, reports no missing character.
\begin{equation} \label{eq:ls-ridge}
    L(\beta) = \frac{1}{2n}\|y - X\beta\|_2^2 + \frac{\mu}{2}\|\beta\|^2_2.
\end{equation}

\begin{assumption} \label{assumption:noise}
    The coordinates of $\varepsilon$ satisfy $\mathbb{E}[\varepsilon_i \mid x_i] = 0$ and $\mathbb{E}[\varepsilon_i^2 \mid x_i] = \tau^2 < \infty$ for all $1 \leq i \leq n$.
\end{assumption}

\begin{corollary}[Ridgeless risk] \label{cor:gen-risk-betak}
Let $X = U \Sigma_X V^T$ be any feature matrix and $y = X \beta_* + \varepsilon$. Assume that Assumption~\ref{assumption:noise} holds. If $\beta_k$ are the parameters after $k$ steps of gradient descent for the least squares problem~\eqref{eq:ls-ridge} with $\mu = 0$, initialized at any $\beta_0$, and with arbitrary learning rate schedule $\{ \eta_k\}_{k \geq 1}$, then, recalling that $\Phi(k) \equiv \Phi(k; 0)$,
\[
    R(\beta_k) = \left\|\Sigma^{1/2}V\Phi(k)(\tilde{\beta}_0 - \tilde{\beta}_*)\right\|_2^2 + \tau^2\left\|\Sigma^{1/2}V(I - \Phi(k))\Sigma_X^\dagger \right\|_F^2.
\]
\end{corollary}

\begin{assumption} \label{assumption:differentiable-extension}
For all fixed $\zeta$ and learning rate schedules $\{\eta_i\}_{i \geq 1}$ such that for all $i$, $\eta_i \le \zeta^{-1}$, the function $k \mapsto \phi(k; \zeta, \{\eta_i\})$ can be extended to a monotonic differentiable function on $[1, \infty)$.
\end{assumption}

\begin{assumption} \label{assumption:spherical}
    The entries of $\beta_0 - \beta_*$ are i.i.d.\ and have mean 0 and variance $\sigma^2$.
\end{assumption}

\begin{theorem}[Early stopping] \label{thm:earlystopping}
Let $X = U \Sigma_X V^T$ be any feature matrix with $\mathrm{rank}(X) = r$, and $y = X \beta_* + \varepsilon$. Recall that $\Lambda_1$ is the largest eigenvalue of $X^T X$. Suppose that Assumptions~\ref{assumption:noise}, \ref{assumption:differentiable-extension}, and~\ref{assumption:spherical} hold.
Let $\beta_k$ be the parameters after $k$ steps of gradient descent for the unpenalized least squares problem~\eqref{eq:ls-ridge} with $\mu = 0$, initialized at any $\beta_0$.
If the learning rate schedule $\{ \eta_k \}_{k \geq 1}$ is such that $\eta_k \le 1 / (n^{-1} \Lambda_1)$ for all $k$, and for all $j = 1, \ldots, r$ we have that 
\begin{equation} \label{eq:earlystopping_condition}
    \lim_{k \to \infty} \frac{\phi(k; n^{-1} \Lambda_j)}{\phi(0; n^{-1} \Lambda_j)} < \frac{\tau^2}{\Lambda_j \sigma^2 + \tau^2},
\end{equation}
then there is a finite $T$ such that for all $k \ge T$, $R_B(\beta_k) \ge R_B(\beta_T)$. That is, early stopping is beneficial.
Furthermore, if $P := V^T\Sigma V$, then for all $k$, 
\begin{equation} \label{eq:earlystopping_lower}
   R_B(\beta_k) \ge \sigma^2\left[ \sum_{j=1}^r P_{jj}\frac{\tau^2}{\Lambda_j\sigma^2 + \tau^2} + \sum_{j=r+1}^p P_{jj}\right].
\end{equation}
\end{theorem}

\begin{proof}[Proof.]
From Corollary~\ref{cor:gen-risk-betak}, noting that $\Lambda^\dagger = (\Sigma_X^T \Sigma_X)^{\dagger}$ is a rank $r$ diagonal matrix, the expected excess risk after $k$ iterations is given by
\begin{align*} 
    R(\beta_k) &= (\tilde{\beta}_0 - \tilde{\beta}_*)^T \Phi(k) P \Phi(k) (\tilde{\beta}_0 - \tilde{\beta}_*) + \tau^2\Tr\left( \Lambda^\dagger (I-\Phi(k))P(I-\Phi(k)) \right)\\
    &= \sum_{i,j=1}^p P_{ij} \left[\frac{\phi(k; n^{-1}\Lambda_i)\phi(k; n^{-1}\Lambda_j)}{\phi(0; n^{-1}\Lambda_i)\phi(0; n^{-1}\Lambda_j)}(\tilde{\beta}_0 - \tilde{\beta}_*)_i(\tilde{\beta}_0 - \tilde{\beta}_*)_j \right] \\
    &\quad+ \sum_{j=1}^r P_{jj} \cdot \tau^2\frac{1}{\Lambda_j}\left(1-\frac{\phi(k; n^{-1}\Lambda_j)}{\phi(0; n^{-1}\Lambda_j)}\right)^2.
\end{align*}
Taking the expectation over $\beta_0 - \beta_*$ under Assumption~\ref{assumption:spherical}, we see that the cross terms vanish, and the Bayes excess risk is given by
\begin{equation} \label{eq:earlystopping_1}
    R_B(\beta_k) = \sum_{j=1}^{p} P_{jj} \cdot \frac{\phi(k;n^{-1}\Lambda_j)^2}{\phi(0;n^{-1}\Lambda_j)^2}\sigma^2 
    + \sum_{j=1}^r P_{jj} \cdot \tau^2\frac{1}{\Lambda_j}\left(1-\frac{\phi(k; n^{-1}\Lambda_j)}{\phi(0; n^{-1}\Lambda_j)}\right)^2.
\end{equation}
Taking the derivative in $k$ and noting that $\phi(k; n^{-1}\Lambda_j)$ is a constant for $j > r$, we get
\begin{align}
    \partial_k R_B(\beta_k)
    &= 2\sum_{j=1}^r P_{jj} \cdot \frac{\partial_k \phi(k; n^{-1}\Lambda_j)}{\phi(0; n^{-1}\Lambda_j)} \left[\frac{\phi(k; n^{-1}\Lambda_j)}{\phi(0; n^{-1}\Lambda_j)}\sigma^2 - \tau^2 \frac{1}{\Lambda_j}\left(1-\frac{\phi(k; n^{-1}\Lambda_j)}{\phi(0;n^{-1}\Lambda_j)}\right) \right] \nonumber\\
    &= 2\sum_{j=1}^r \frac{P_{jj}}{\Lambda_j} \cdot \frac{\partial_k \phi(k; n^{-1}\Lambda_j)}{\phi(0; n^{-1}\Lambda_j)} \left[\frac{\phi(k;n^{-1}\Lambda_j)}{\phi(0; n^{-1}\Lambda_j)} \Big(\sigma^2 \Lambda_j + \tau^2\Big) - \tau^2  \right]. \label{eq:risk-derivative-mu0}
\end{align}
Note that since $\eta_k \le 1/(n^{-1} \Lambda_1)$ for all $k$, $\phi(k; n^{-1}\Lambda_j) / \phi(0; n^{-1} \Lambda_j)$ is a non-increasing function of $k$ by Assumption~\ref{assumption:differentiable-extension}, and hence $\partial_k \phi(k; n^{-1}\Lambda_j) / \phi(0; n^{-1} \Lambda_j) \leq 0$. Since $\Lambda_j \ge 0$ and $P_{jj} \geq 0$ (because $P$ is a positive semidefinite matrix), the sign of the derivative $\partial_k R_B(\beta_k)$ is determined by 
\[
    \frac{\phi(k;n^{-1}\Lambda_j)}{\phi(0; n^{-1}\Lambda_j)} \Big(\sigma^2 \Lambda_j + \tau^2\Big) - \tau^2.
\]
Therefore, by using the assumption~\eqref{eq:earlystopping_condition} that 
\[
    \lim_{k \to \infty} \frac{\phi(k;n^{-1}\Lambda_j)}{\phi(0; n^{-1}\Lambda_j)} < \frac{\tau^2}{\sigma^2\Lambda_j + \tau^2},
\]
we have that for sufficiently large $k$,
\[
    \partial_k R_B(\beta_k) \ge 2 \sum_{j=1}^{r} \frac{P_{jj}}{\Lambda_j} \cdot \frac{\partial_k \phi(k; n^{-1} \Lambda_j)}{\phi(0; n^{-1} \Lambda_j)} (\tau^2 - \tau^2) = 0.
\]
This shows that the derivative of the expected excess risk is eventually positive, from which we conclude that we should have stopped earlier to minimize the risk.

To obtain the lower bound on the Bayes excess risk, observe that we can minimize the expected excess risk independently in each eigendirection. Since $\phi$ is monotonic, we see that solving 
\[
   \frac{\phi(k_j;n^{-1}\Lambda_j)}{\phi(0; n^{-1}\Lambda_j)} = \frac{\tau^2}{\sigma^2\Lambda_j + \tau^2}
\]
for $k_j$ achieves a global minimum in the $j$th eigendirection for $j \le r$. By plugging this into the expression~\eqref{eq:earlystopping_1} for the Bayes excess risk for each $j \le r$, we obtain the lower bound as follows. For $j \le r$, the corresponding term in the summation is $P_{jj}$ times 
\begin{align*}
    \sigma^2 \left(\frac{\tau^2}{\sigma^2\Lambda_j + \tau^2}\right)^2 + \tau^2\frac{1}{\Lambda_j}\left(1-\frac{\tau^2}{\sigma^2\Lambda_j + \tau^2}\right)^2
    &= \sigma^2 \left(\frac{\tau^2}{\sigma^2\Lambda_j + \tau^2}\right)^2 + \tau^2\frac{1}{\Lambda_j}\left(\frac{\sigma^2\Lambda_j}{\sigma^2\Lambda_j + \tau^2}\right)^2 \\
    &= \frac{\sigma^2 \tau^2}{\sigma^2\Lambda_j + \tau^2}.
\end{align*}
For $j > r$, note that $\phi(k;n^{-1}\Lambda_j) / \phi(0; n^{-1}\Lambda_j) = 1$. This completes the proof.
\end{proof}

\end{document}
